% =====================================================================
% DRAFT: §1 Introduction and §2 Definitions in the positive frame.
%
% Self-contained excerpt — drop into main.tex by replacing the existing
% \section{Introduction} ... up through the end of \section{Definitions}.
% Macro definitions (\dotop, \core, \FRM, \ICP, \Fin) are unchanged
% from main.tex.
%
% Substantive changes vs. the current §1, §2:
%   - The setting is introduced as a positive primitive (finite extensional
%     algebra with binary judgment), not as the residual of four obstructions.
%   - The three capabilities S, D, C are the natural questions about that
%     primitive, motivated before being defined.
%   - K-infinity moves out of "what narrows the candidate class" into a
%     contextualising paragraph relating this work to the PCA/TCA tradition.
%   - Single-absorber collapse demotes from Proposition 1 to a remark
%     (the cardinality 2 is part of the definition of "binary judgment",
%     so the collapse just records why "binary" is the minimum).
%   - The "tour of obstructions" paragraph in §1 becomes a single
%     positive thread: the no-associativity / no-mediality / no-RSD /
%     no-LSD / no-right-identity cluster is presented as a single
%     classifier-constancy schema with a unified motivation.
%   - Asymmetry (Theorem 5 in current numbering) is mentioned in §1
%     and shown to follow from the absorber semantics directly, not
%     listed as an obstruction.
%
% No theorem statement, no proof, and no Lean artifact moves. Section 3
% (Universal Theorems) and beyond are untouched.
% =====================================================================

\section{Introduction}
\label{sec:introduction}

A finite magma is the Cayley table of a binary operation, stripped of
syntax and types. We study a particular kind: \emph{finite extensional
algebras with binary judgment} --- finite magmas whose binary operation
returns, alongside data outputs, judgments drawn from a fixed two-element
subset of the carrier playing the role of truth values. Within this
setting we ask three natural operational questions: can the algebra
\emph{represent} its own data, can it \emph{classify} its own structure,
and can it \emph{compose} its own operations. Each question has a precise
formal answer (S, D, C), and our main result is that the three are
pairwise independent: no one implies any other.

The setting is fixed by three structural commitments, all of them positive.
\emph{Extensionality} is the foundational one: an element's identity is
exhausted by its row in the Cayley table, so an element \emph{is} its
action under the operation~\cite{burris81}. \emph{Two distinguished
truth values} $z_1, z_2$ are the minimum number for the operation to
register a binary outcome; equipping the algebra with judgment cardinality
one is degenerate (Remark~\ref{rem:single-judgment}), and higher arities
are natural generalisations we do not pursue. \emph{Inertness of the
truth values} formalises what makes them truth values rather than data:
$z_i \dotop x = z_i$ for all $x$, so post-judgment behaviour depends only
on judgment status. The truth values are operationally distinguishable
(they absorb), structurally interchangeable (no axiom prefers $z_1$ to
$z_2$), and exhaustive among inert elements (no third element is
left-absorbing). Together these commitments select the class of finite
extensional 2-pointed magmas.

The three capabilities. Let $\core(S) = S \setminus \{z_1, z_2\}$ denote
the carrier minus its truth values. \emph{Internal splitting~(S)}: there
exists a section/retraction pair $(s, r) \in \core(S)^2$ with
$r \dotop (s \dotop x) = x$ on $\core(S)$, so the algebra has internal
encoding/decoding for its own data. \emph{Internal dichotomy~(D)}: every
element of core has the same judgment behaviour across all of core ---
either every output lies in $\{z_1, z_2\}$ (a \emph{classifier}) or no
output does (a \emph{non-classifier}) --- and at least one example of each
type exists. \emph{Internal composition~(C)}: there exist three pairwise
distinct core elements $a, b, c$ with $b$ core-preserving and
$a \dotop x = c \dotop (b \dotop x)$ on $\core(S)$, with $a$ taking at
least two distinct values --- a non-trivial factorization in the
left-regular representation. Each capability is what an external machine
would otherwise have to supply (naming and decoding, binary judgment,
operation chaining), realised instead as a behaviour of elements of the
algebra modulo the operation.

\paragraph{Three structural results.} Pairwise independence (Section~\ref{sec:independence}):
no one of the three capabilities implies any other. Lean-verified finite
counterexamples at sizes $N \in \{4, 5, 6, 8, 10\}$ establish all six
non-implications; four are tight. Joint irredundance
(Section~\ref{sec:boolean-cube}): all eight Boolean combinations of
$(S, D, C)$ are realised by Lean-verified witnesses, so no capability is
a Boolean function of the other two. Optimal coexistence
(Section~\ref{sec:bound}): the smallest cardinality at which all three
capabilities can hold simultaneously is $N = 5$, and this is tight: ICP
requires three pairwise-distinct core elements, forcing $|S| \geq 5$.
Each Cayley table is found by SAT search and independently verified in
Lean~4~\cite{moura21} with zero \texttt{sorry}.

\paragraph{What follows from binary judgment.} Several classical
algebraic identities are incompatible with this setting. The proofs
share a single schema: any equation whose specialisation at a
classifier $\tau$ forces $\tau$'s row to be constant collides with
extensionality, equating $\tau$ with one of the truth values. From D
alone: commutativity, mediality, right self-distributivity all fail
(Theorems~\ref{thm:asymmetry}, \ref{thm:no-mediality}, \ref{thm:no-rsd}).
From S+D: associativity, left self-distributivity, right identity all
fail (Theorems~\ref{thm:no-assoc}, \ref{thm:no-lsd}, \ref{thm:no-rid}).
The classical algebraic hierarchy --- semigroups, monoids, groups, rings
--- therefore cannot host all three capabilities, despite each capability
being individually expressible in classical terms. The intuition behind
the schema is that binary judgment is intrinsically non-chainable:
composing the operation through a classifier collapses the judgment to
a fixed verdict, eliminating discriminative content.

\paragraph{What lies beyond.} Two structural results extend the picture.
The three-category decomposition $S = Z \sqcup C \sqcup N$ induced by D
is an isomorphism invariant (Theorem~\ref{thm:functoriality}); the
unique coarsest non-trivial isomorphism-invariant partition of core is
the classifier/non-classifier split (Corollary~\ref{cor:canonicity}).
At $N = 5$ a structure theorem (Theorem~\ref{thm:N5-structure}) pins the
role assignments of retraction, classifier, and ICP-triple onto the
three core slots with no slack: every retraction pair has $s = r$, every
ICP-triple has the form $(\tau_1, g, \tau_2)$, and the only possible
non-trivial automorphism is the classifier transposition
$(\tau_1\,\tau_2)$ (Theorem~\ref{thm:mirror-row}). At $N \geq 6$ strong S
re-opens and the role-shape fractures (Proposition~\ref{prop:N6-nonrigid}),
but partial invariants survive across the transition
(Propositions~\ref{prop:N6-aut-structure}, \ref{prop:N6-h-triples}).

\paragraph{Relation to universality programmes.} Combinatory algebras
(PCAs, TCAs)~\cite{vanoosten08, longley15, bethke88} and Turing
categories~\cite{cockett08} solve a different problem: full internal
universality of the operation, paid for by infinite cardinality
(Theorem~\ref{thm:k-infinity}, the K-infinity result of the PCA tradition).
The setting of this paper retains finite cardinality and obtains, in
exchange, a weaker form of universality --- representation without a
global projector, classification without a subobject classifier,
composition without a total internal hom. The three capabilities S, D, C
are what survives the trade. Section~\ref{sec:related} discusses the
finite-vs-infinite cardinality dichotomy in more detail; it is context
for our setting, not a fact derivable inside it.

\paragraph{Contributions.}
\begin{enumerate}
  \item A positive axiomatisation of the setting as finite extensional
    algebras with binary judgment, together with a uniform
    classifier-constancy schema explaining why associativity, mediality,
    self-distributivity, right identity, and commutativity all fail
    (Section~\ref{sec:universal}).
  \item An independence theorem: S, D, and C are pairwise independent,
    with Lean-verified counterexamples in all six directions at $N = 4$
    through $N = 10$ (Section~\ref{sec:independence}), and coexistence
    witnesses at $N = 5$ and $N = 6$ (Section~\ref{sec:bound}).
  \item An optimal bound: $N = 5$ is the minimum carrier size for S+D+C
    coexistence; ICP is impossible at $N = 4$ by pigeonhole
    (Theorem~\ref{thm:no-icp-4}).
  \item Decomposition invariance and role rigidity: the $Z/C/N$ partition
    is preserved by isomorphisms (Section~\ref{sec:invariance}); the
    $N \geq 6$ principal coexistence witness and the smaller S+D witnesses
    have trivial automorphism group, while the canonical $N = 5$ S+D+C
    witness internalises a classifier-swap as the action of its
    non-classifier (Proposition~\ref{prop:witness-rigid},
    Remark~\ref{rem:canonical-witness}); at $N = 5$ a structure theorem
    forces retraction, classifier, and C-triple onto the three core slots
    and rules out strong S, and the mirror-row theorem
    (Theorem~\ref{thm:mirror-row}) bounds $|\mathrm{Aut}(M)| \leq 2$;
    at $N \geq 6$ strong S returns and rigidity fails
    (Proposition~\ref{prop:N6-nonrigid}).
  \item A characterization of C: the ICP is logically equivalent to the
    standard Compose+Inert axioms under variable renaming, with
    Lean-verified necessity of the non-degeneracy guard conditions
    (Section~\ref{sec:icp}).
  \item Joint irredundance: all eight cells of the $(S, D, C)$ Boolean
    cube are inhabited by Lean-verified witnesses
    (Section~\ref{sec:boolean-cube}), so no capability is a Boolean
    function of the other two.
\end{enumerate}

% =====================================================================

\section{Definitions}
\label{sec:definitions}

Throughout, $(S, \dotop)$ is a finite magma on
$\Fin(n) = \{0, 1, \ldots, n-1\}$ with binary operation
$\dotop : \Fin(n) \times \Fin(n) \to \Fin(n)$.

\begin{definition}[Finite Extensional Algebra with Binary Judgment]
\label{def:epm}
A \emph{finite extensional algebra with binary judgment} --- equivalently,
an \emph{extensional 2-pointed magma} --- is a tuple
$(S, \dotop, z_1, z_2)$ where:
\begin{enumerate}
  \item $z_1, z_2 \in S$ are distinct elements designated as the
    \emph{truth values} of the algebra.
  \item The truth values are \emph{operationally inert}: $z_i \dotop x = z_i$
    for all $x \in S$. Equivalently, $z_1$ and $z_2$ are left-absorbers.
  \item No other element is operationally inert: no other element is a
    left-absorber.
  \item \emph{Extensionality}: distinct elements have distinct rows in
    the Cayley table.
\end{enumerate}
We write $\core(S) := S \setminus \{z_1, z_2\}$ for the elements of $S$
that are not truth values.
\end{definition}

The four clauses isolate four structural commitments. Clause~(i) declares
two distinguished judgment outputs and requires them to be distinct ---
``binary'' means two-valued. Clause~(ii) gives the truth values their
operational meaning: their post-judgment behaviour depends only on
their judgment status, not on what they are judging. Clause~(iii) closes
the truth-value subset against operational impostors: every other element
of $S$ varies under left-multiplication, so no other element ``acts like
a truth value'' without being one. Clause~(iv) is the foundational
extensional commitment: an element's row in the Cayley table is a
complete operational description of it, so equality of behaviour is
equality.

The choice of cardinality two for the judgment subset is forced by
non-degeneracy:

\begin{remark}[Why binary, not unary]
\label{rem:single-judgment}
If a candidate algebra had only one truth value $z$, any element
$\tau$ purportedly classifying inputs against $\{z\}$ would satisfy
$\tau \dotop x = z$ for all $x$. The row of $\tau$ would then equal the
row of $z$ (both constantly $z$), and extensionality would force
$\tau = z$ --- the supposed classifier collapses into the truth value
itself. Binary judgment is therefore the smallest non-degenerate case,
not a parameter choice; $k$-valued generalisations for $k \geq 3$
strengthen the theory in a different direction we do not pursue here.
Lean: \texttt{single\_absorber\_collapse}.
\end{remark}

The clause~(iii) requirement that no other element be a left-absorber
is independent of clauses~(i), (ii), and (iv) together with the
classifier dichotomy of Definition~\ref{def:cd}: at every $N \in \{4, 5, 6\}$,
Z3 produces an extensional 2-pointed magma satisfying D in which a third
core element has a constant row (script
\texttt{scripts/canonicality/no\_other\_zeros\_independence.py}). The
clause is therefore a genuine axiom, not a consequence of the rest.

We do \emph{not} commit to a stronger reading of the two truth values as
``opposites'' in any structural sense. The theorems are compatible both
with a symmetric reading (two indistinguishable judgment outputs) and a
polar reading (two poles of a polarity). Whether $z_1$ corresponds to
``false'' and $z_2$ to ``true'' or vice versa is a free choice the
axioms do not constrain. This freedom is genuine: applications can
specialise the reading without changing the theorems.

\begin{definition}[Faithful Retract Magma]
\label{def:frm}
A \emph{Faithful Retract Magma} ($\FRM$) extends an extensional 2-pointed
magma with a \emph{retraction pair} $s, r \in \core(S)$ satisfying
$r \dotop (s \dotop x) = x$ for all $x \in \core(S)$, with
$r \dotop z_1 = z_1$ (anchoring). The remaining absorber outputs of $s$
and $r$ are unconstrained.
\end{definition}

\paragraph{Retraction pair convention.}
\label{rem:mutual}
Throughout this paper the retraction pair $(s, r)$ satisfies the
mutual-inverse condition on core: $r \dotop (s \dotop x) = x$ and
$s \dotop (r \dotop x) = x$ for all $x \in \core(S)$. This matches the
Lean formalisation. The choice is load-bearing, not stylistic.
One-sided retraction ($r \dotop (s \dotop x) = x$) is a
decodable-encoding condition: $s$ produces codes that $r$ can recover.
Mutual inverse is the stronger, faithful-representation condition: $s$
and $r$ invert each other in both directions on core, so the encoding
neither invents nor loses information. Most theorems (9 of 11) go
through under the one-sided hypothesis, and we note this where relevant;
two results, retraction pair placement (Theorem~\ref{thm:placement}) and
the tight bound $|S| \geq 5$ when $s \neq r$ (Theorem~\ref{thm:card}),
genuinely require the full mutual inverse.
Proposition~\ref{prop:one-sided-sep} shows this distinction is
non-vacuous.

\begin{proposition}[One-Sided / Mutual-Inverse Separation]
\label{prop:one-sided-sep}
There is an extensional 2-pointed magma on $\Fin(4)$ admitting a
one-sided retraction pair $(s, r)$ with $s \neq r$, $s, r \in \core$,
$r \dotop z_1 = z_1$, and $r \dotop (s \dotop x) = x$ for all
$x \in \core$, but no mutual-inverse pair $(s', r')$ with $s' \neq r'$.
The explicit Cayley table (rows indexed by the left operand) is:
\[
\begin{array}{c|cccc}
\dotop & 0 & 1 & 2 & 3 \\ \hline
0 & 0 & 0 & 0 & 0 \\
1 & 1 & 1 & 1 & 1 \\
2 & 1 & 0 & 1 & 3 \\
3 & 0 & 2 & 1 & 3
\end{array}
\]
with $z_1 = 0$, $z_2 = 1$, $s = 2$, $r = 3$. Mutual inverse fails at
$x = 2$: $s \dotop (r \dotop 2) = 2 \dotop 1 = 0 \neq 2$. Lean:
\texttt{mutual\_strictly\_stronger\_than\_one\_sided}.
\end{proposition}

\begin{proof}
All axioms are verified by direct inspection; Lean discharges each by
\texttt{decide}, including the non-existence of a mutual-inverse pair
with $s' \neq r'$ over $\Fin(4) \times \Fin(4)$. The separation
establishes that mutual inverse strictly implies one-sided retraction,
so the choice of axiom is substantive, and that the tight bound of
Theorem~\ref{thm:card} cannot be weakened to one-sided.
\end{proof}

The retraction pair gives the algebra a section/retraction structure on
core: $s$ acts as a section and $r$ as a retraction. We do not assume a
canonical global encoding.

\begin{definition}[Classifier Dichotomy]
\label{def:cd}
An extensional 2-pointed magma satisfies the \emph{classifier dichotomy}
when:
\begin{enumerate}
  \item A \emph{classifier} $\tau \in \core(S)$ exists with
    $\tau \dotop x \in \{z_1, z_2\}$ for all $x \in S$.
  \item For every $y \in \core(S)$, either
    \begin{itemize}
      \item $y \dotop x \in \{z_1, z_2\}$ for all $x \in \core(S)$ ---
        $y$ is a \emph{classifier}, or
      \item $y \dotop x \notin \{z_1, z_2\}$ for all $x \in \core(S)$ ---
        $y$ is a \emph{non-classifier}.
    \end{itemize}
  \item At least one non-classifier exists (the dichotomy is not
    one-sided).
\end{enumerate}
\end{definition}

The three clauses split into structural and non-vacuity content.
Clauses~(i) and~(ii) together, $D_{\mathrm{struct}}$, are the structural
fragment: the algebra contains \emph{some} classifier and the
classification is global. Clause~(iii) ensures both classes are
inhabited. In the joint S+D+C setting, clause~(iii) is automatic: the
middle element $b$ of any ICP triple is core-preserving (Definition~\ref{def:icp}
clause~1), hence a non-classifier, hence supplies the missing witness
(Lean: \texttt{DStructRetractMagma.has\_non\_classifier\_of\_icp}).
Theorem~\ref{thm:d-factorization} below makes the split formal:
$D \iff D_{\mathrm{struct}} \wedge P$, where $P$ asserts the existence
of a core-preserving element.

The dichotomy creates a three-category decomposition
$S = Z \sqcup C \sqcup N$ where $Z = \{z_1, z_2\}$ (truth values),
$C$ (classifiers), and $N$ (non-classifiers). No element can be partially
classifying. We call outputs in $\{z_1, z_2\}$ \emph{absorber-valued}.

\begin{definition}[Internal Composition Property]
\label{def:icp}
A magma $(S, \dotop, z_1, z_2)$ satisfies the \emph{Internal Composition
Property} ($\ICP$) when:
\[
\exists\, a, b, c \in \core(S),\ \text{pairwise distinct, such that:}
\]
\begin{enumerate}
  \item \emph{Core-preservation}: $b \dotop x \notin \{z_1, z_2\}$ for
    all $x \in \core(S)$.
  \item \emph{Factorization}: $a \dotop x = c \dotop (b \dotop x)$ for
    all $x \in \core(S)$.
  \item \emph{Non-triviality}: $|\{a \dotop x : x \in \core(S)\}| \geq 2$.
\end{enumerate}
\end{definition}

ICP says that the left-regular representation contains a non-trivial
factorization on core: one core action factors through two others, one
of them core-preserving --- a single fragment of internal composition.
The pairwise-distinctness and non-triviality clauses are guard conditions
blocking degeneracies; both are proved necessary by Lean-verified
witnesses (Section~\ref{sec:icp}). In Section~\ref{sec:icp} we prove
that $\ICP$ is equivalent to the standard operational axioms
Compose+Inert (Theorem~\ref{thm:icp-equiv}).

\begin{definition}[The Three Capabilities]
\label{def:capabilities}
An extensional 2-pointed magma has:
\begin{itemize}
  \item Capability $S$ (\emph{internal splitting}) if it admits a
    retraction pair (Definition~\ref{def:frm}).
  \item Capability $D$ (\emph{internal dichotomy}) if it satisfies the
    classifier dichotomy (Definition~\ref{def:cd}).
  \item Capability $C$ (\emph{internal composition}) if it satisfies the
    Internal Composition Property (Definition~\ref{def:icp}).
\end{itemize}
\end{definition}

All three capabilities are properties of the same base structure. None
is defined as an extension of any other, making the independence theorem
a statement about six non-implications, all proved by Lean-verified
counterexamples.

\paragraph{Why these three.} Each capability answers a specific positive
question about a finite extensional algebra with binary judgment, in a
way that is operational rather than typed. Extensionality identifies
each element with its row in the Cayley table, so the elements of the
algebra are a family of typed maps
$\{L_a^{\core} : \core \to \core \cup Z\}$, one per element, with the
truth-value maps $L_{z_1}, L_{z_2}$ being constant. S, D, and C
correspond to three fundamental structural questions about such a
family:
\begin{itemize}
  \item \emph{Splitting} (S): does the identity on core split through
    the family --- do some $L_s, L_r$ satisfy
    $L_r \circ L_s = \mathrm{id}_{\core}$? This is the standard
    section-retraction pair~\cite{maclane71}.
  \item \emph{Codomain purity} (D): does each map hit $\core$ or $Z$,
    never both? This is the unique coarsest non-trivial
    isomorphism-invariant partition of core
    (Corollary~\ref{cor:canonicity}).
  \item \emph{Factorization} (C): does the family contain a non-trivial
    composition $L_a = L_c \circ L_b$ with $L_b$ core-preserving and
    $L_a$ non-constant? This is the weakest factorization condition
    that is not vacuously satisfiable (Lean-verified separations in
    \texttt{ICP.lean}).
\end{itemize}
Each property is minimal in its own sense. S is the standard notion of
splitting. D is the coarsest invariant partition. C is the weakest
non-degenerate factorization. The independence and joint-irredundance
results (Theorems~\ref{thm:independence}, \ref{thm:cube}) confirm these
are three questions, jointly irredundant.

\paragraph{Other structural questions.} The slice $(S, D, C)$ is one
choice within a broader axiom landscape. A Z3 canonicality probe
(\texttt{scripts/canonicality/probe\_axioms.py}) tested thirteen
candidate operational axioms on extensional 2-pointed magmas at
$N \in \{5, 6\}$ via pairwise directional implication queries against
$S$, $D$, $C$; ten are pairwise independent of $(S, D, C)$ in both
directions at both sizes, three collapse into derived consequences
(retraction pair gives one-sided retraction and left-cancellation; both
$D$ and $C$ imply $P$ via Theorem~\ref{thm:d-factorization}). $(S, D, C)$
is therefore one structurally selectable slice among many. What
distinguishes it is that the canonical $N = 5$ S+D+C witness is fully
determined by the topos-theoretic indicator reading of $D$ together with
the requirement that the magma's automorphism be internalised in its
operation (Remark~\ref{rem:canonical-witness}); the other ten probe
axioms have no such joint structural characterisation.

\paragraph{The internalisation frame.}
The magma operation $\dotop$ is taken as given throughout --- the
structuring relation of $(A, \dotop)$, not itself an object inside the
structure. This is not a gap but a structural commitment: no finite
algebra can contain its own operation as an element without generating
a meta-operation, a Tarski/Lawvere-style regress. We study what
\emph{can} be internalised within $(A, \dotop)$ modulo $\dotop$:
services that would otherwise be supplied externally (naming and
decoding, binary judgment, operation chaining), realised instead as
behaviours of elements relative to the fixed operation. The three
capabilities S, D, C are each one such internalisation; the operation
$\dotop$ is the one component deliberately left outside the
internalisation catalog.
